\documentclass[12pt,a4paper]{article}

% --------------------------------------------------
% Pacotes
% --------------------------------------------------
\usepackage[utf8]{inputenc}
\usepackage[T1]{fontenc}
\usepackage[brazil]{babel}

\usepackage{amsmath}
\usepackage{amssymb}
\usepackage{graphicx}
\usepackage{booktabs}
\usepackage{float}
\usepackage{geometry}
\usepackage{setspace}
\usepackage{indentfirst}
\usepackage{hyperref}

% --------------------------------------------------
% Configurações
% --------------------------------------------------
\geometry{
    a4paper,
    left=3cm,
    right=2cm,
    top=3cm,
    bottom=2cm
}

\onehalfspacing

% --------------------------------------------------
% Informações do documento
% --------------------------------------------------
\title{ME524 -- Laboratório 1}
\author{Seu Nome}
\date{\today}

% --------------------------------------------------
% Documento
% --------------------------------------------------
\begin{document}

\maketitle

\tableofcontents
\newpage

% --------------------------------------------------
\section{Introdução}
% --------------------------------------------------

Este documento apresenta as análises realizadas no Laboratório 1
da disciplina ME524.

% --------------------------------------------------
\section{Metodologia}
% --------------------------------------------------

Nesta seção será apresentada a metodologia utilizada na análise.

% --------------------------------------------------
\section{Análise Exploratória}
% --------------------------------------------------

Nesta seção serão apresentados os resultados da análise exploratória
dos dados.

% Exemplo de equação:
\begin{equation}
    \bar{x} = \frac{1}{n}\sum_{i=1}^{n}x_i
\end{equation}

% --------------------------------------------------
\section{Resultados}
% --------------------------------------------------

Os principais resultados obtidos são apresentados nesta seção.

% Exemplo de figura:
%
% \begin{figure}[H]
%     \centering
%     \includegraphics[width=0.8\textwidth]{figures/exemplo.png}
%     \caption{Exemplo de figura.}
%     \label{fig:exemplo}
% \end{figure}

% Exemplo de tabela:
\begin{table}[H]
    \centering
    \caption{Exemplo de tabela.}
    \begin{tabular}{lcc}
        \toprule
        Variável & Média & Desvio-padrão \\
        \midrule
        $X_1$ & 10.5 & 2.1 \\
        $X_2$ & 15.3 & 3.4 \\
        \bottomrule
    \end{tabular}
    \label{tab:exemplo}
\end{table}

% --------------------------------------------------
\section{Conclusão}
% --------------------------------------------------

Nesta seção serão apresentadas as principais conclusões da análise.

% --------------------------------------------------
% Referências
% --------------------------------------------------

% Caso utilize um arquivo .bib:
%
% \bibliographystyle{plain}
% \bibliography{references}

\end{document}